% ============================================================
% secciones/03_metodologia.tex
% ============================================================

\section{Metodología}
\label{sec:metodo}

\subsection{Diseño del estudio}

Se utiliza un diseño de validación cruzada estratificada de $k=5$ pliegues, estratificando por la variable respuesta $N$ para preservar la proporción de pólizas sin siniestros en cada partición. El periodo 2019–2022 se utiliza para entrenamiento y el año 2023 como muestra de prueba externa (hold-out temporal).

\subsection{Métricas de evaluación}

\textbf{Frecuencia:}
\begin{equation}
  \text{RMSE}_N = \sqrt{\frac{1}{n}\sum_{i=1}^{n}(\hat{N}_i - N_i)^2}
\end{equation}

\textbf{Severidad} (sobre pólizas con al menos un siniestro):
\begin{equation}
  \text{RMSE}_X = \sqrt{\frac{1}{m}\sum_{j=1}^{m}(\hat{X}_j - X_j)^2}
\end{equation}

\textbf{Prima pura:}
\begin{equation}
  \text{RMSE}_\Pi = \sqrt{\frac{1}{n}\sum_{i=1}^{n}(\hat{\Pi}_i - \Pi_i)^2}, \quad \hat{\Pi}_i = \hat{N}_i \cdot \hat{X}_i
\end{equation}

Se complementa con el AIC y BIC para los modelos GLM, y la ganancia de Gini para comparación global de capacidad discriminante \citep{Frees2014}.

\subsection{Selección de variables}

Para los modelos GLM se aplica selección por pasos hacia atrás (backward stepwise) minimizando el AIC. Para XGBoost se utiliza la importancia SHAP (\textit{SHapley Additive exPlanations}) para identificar las variables más relevantes y comparar con el GLM.
